% ==============================================================================
%  acc_paper_experiments_additions.tex  —  SINGLE FILE, paste into main paper
% ==============================================================================
% USE: In main .tex, (1) add \graphicspath{{./figures/}} to preamble;
%      (2) copy all outputs/fig_*.pdf to ./figures/;
%      (3) REPLACE your existing Section V-E tail (Table III + runtime text)
%          AND insert Section V-F with:
%            % ==============================================================================
%  acc_paper_experiments_additions.tex  —  SINGLE FILE, paste into main paper
% ==============================================================================
% USE: In main .tex, (1) add \graphicspath{{./figures/}} to preamble;
%      (2) copy all outputs/fig_*.pdf to ./figures/;
%      (3) REPLACE your existing Section V-E tail (Table III + runtime text)
%          AND insert Section V-F with:
%            % ==============================================================================
%  acc_paper_experiments_additions.tex  —  SINGLE FILE, paste into main paper
% ==============================================================================
% USE: In main .tex, (1) add \graphicspath{{./figures/}} to preamble;
%      (2) copy all outputs/fig_*.pdf to ./figures/;
%      (3) REPLACE your existing Section V-E tail (Table III + runtime text)
%          AND insert Section V-F with:
%            % ==============================================================================
%  acc_paper_experiments_additions.tex  —  SINGLE FILE, paste into main paper
% ==============================================================================
% USE: In main .tex, (1) add \graphicspath{{./figures/}} to preamble;
%      (2) copy all outputs/fig_*.pdf to ./figures/;
%      (3) REPLACE your existing Section V-E tail (Table III + runtime text)
%          AND insert Section V-F with:
%            \input{acc_paper_experiments_additions}
%          immediately BEFORE \section{Conclusion}.
%
% FIGURES (outputs/ → figures/): fig_confidence(Fig1), fig_backoff(Fig2),
%   fig_traj(Fig3), fig_pedestrian(Fig4), fig_intersection(Fig5),
%   fig_montecarlo_combined(Fig6), fig_sensitivity(Fig7, NEW).
% LABELS: tab:montecarlo(Table III), tab:runtime(Table IV, NEW),
%   fig:montecarlo(Fig6), fig:sensitivity(Fig7, NEW), sec:sensitivity(V-F).
% REFRESH numbers: MPLCONFIGDIR=.mplcache python run_all.py --mc-trials 100
%
% FIG 6 (keep in main .tex ABOVE this file; PDF has NO panel titles):
% \begin{figure}[t]
% \centering
% \includegraphics[width=\linewidth]{fig_montecarlo_combined.pdf}
% \caption{(a)~Mean back-off per trial, $100$ moderate-severity trials
% ($0\%$ violations, all methods).
% (b)~Violation rate, all methods, both severity regimes ($100$ trials each).}
% \label{fig:montecarlo}
% \end{figure}
% ==============================================================================

% ----- V-E (replace old Table III + ``Solver runtime'' paragraph) --------------

\begin{table}[t]
\centering
\caption{Monte Carlo results, $100$ trials per cell.}
\label{tab:montecarlo}
\small
\setlength{\tabcolsep}{3pt}
\begin{tabular}{lcccc}
\toprule
 & \multicolumn{2}{c}{Moderate severity} & \multicolumn{2}{c}{Severe (stress test)} \\
Method & B.o.\ [m] & Viol.\ & B.o.\ [m] & Viol.\ \\
\midrule
Fixed $\delta=0.05$ & $4.36\pm0.00$ & $0\%$ & $4.36\pm0.00$ & $20\%$ \\
Worst-case fixed & $9.95\pm0.00$ & $0\%$ & $9.95\pm0.00$ & $10\%$ \\
Naive reactive & $2.45\pm0.24$ & $0\%$ & $2.57\pm0.33$ & $30\%$ \\
Confidence-adaptive & $2.57\pm0.14$ & $0\%$ & $2.62\pm0.16$ & $20\%$ \\
\bottomrule
\end{tabular}
\end{table}

\begin{table}[t]
\centering
\caption{QP solver runtime over all Monte Carlo solves ($T_s=200$\,ms control period). General-purpose CPU; all four methods share identical QP structure.}
\label{tab:runtime}
\small
\begin{tabular}{lcccc}
\toprule
Method & Mean [ms] & Median [ms] & 95th [ms] & Max [ms] \\
\midrule
Fixed $\delta=0.05$  & 3.8 & 3.8 & 4.2  & 38.4  \\
Worst-case fixed     & 4.2 & 3.9 & 5.4  & 111.2 \\
Naive reactive       & 3.9 & 3.8 & 4.9  & 32.8  \\
Confidence-adaptive  & 3.8 & 3.8 & 4.0  & 25.7  \\
\bottomrule
\end{tabular}
\vspace{2pt}
\footnotesize Confidence update overhead is $<0.01$\,ms (closed-form scalar).
\end{table}

\textbf{Moderate regime: efficiency confirmed, no violations.}
Table~\ref{tab:montecarlo} (left) confirms the single-trial result of
Section~\ref{sec:sim-lanechange} in distribution: across $100$ randomized
trials of Scenario~A (braking severity, onset, duration, and measurement noise
randomized per trial), comparing fixed-$\delta=0.05$, worst-case-fixed
($\delta=\delta_{\min}$), naive-reactive (raw $\bar e_k\!\to\!\delta$ without
Proposition~1), and confidence-adaptive, the adaptive controller achieves a
$41.5\%$ mean back-off reduction vs.\ fixed ($2.57$\,m vs.\ $4.36$\,m) with
zero violations, matching naive-reactive's mean but with $41\%$ lower
trial-to-trial variance (Fig.~\ref{fig:montecarlo}a; $\sigma=0.14$\,m vs.\
$0.24$\,m), consistent with the regularized MAP estimator smoothing a purely
reactive rule.

\textbf{Severe regime: an honest trade-off, not a free lunch.}
Widening severity (deceleration $\in[2.0,5.0]$\,m/s$^2$, onset
$\in[3.0,5.0]$\,s, duration $\in[1.0,3.0]$\,s) yields nonzero violations
for all methods (Table~\ref{tab:montecarlo}, right), including worst-case
fixed ($10\%$), indicating some trials exceed what $a_{\max}=2.5$\,m/s$^2$
can compensate; this reflects the scenario distribution, not a particular
risk policy. Margin-reducing methods show higher violation rates (naive
$30\%$, adaptive $25\%$) than conservative ones (fixed $20\%$, worst-case
$10\%$), the expected cost of relaxing margins during calm periods
(Theorem~\ref{thm:feasible}). Adaptive is consistently safer than naive
(lower violation rate and variance), but adaptation is not a substitute for
choosing $\delta_{\min}$, $\delta_{\max}$, $\kappa$ conservatively enough---a
calibration question flagged as future work.

\textbf{Solver runtime.}
Table~\ref{tab:runtime} reports wall-clock QP time pooled over ${\approx}40{,}000$
solves (both regimes, four methods). Mean $3.8$\,ms, median $3.8$\,ms, 95th
$4.0$\,ms---well within $T_s=200$\,ms. Max $111$\,ms occurs on OSQP-fallback
(SCS) instances; a deployed system would use a solve-time watchdog (reuse
previous input or emergency maneuver).

% ----- V-F (NEW: insert before \section{Conclusion}) --------------------------

\subsection{Parameter Sensitivity Analysis}
\label{sec:sensitivity}

Parameters $\alpha$, $\beta$ (Proposition~1), $\lambda$ (Eq.~\eqref{eq:delta-of-c}),
and $\kappa$ (Eq.~\eqref{eq:sigma-eff}) each play a distinct role.
Fig.~\ref{fig:sensitivity} sweeps one parameter at a time on Scenario~A
(adaptive controller), holding others at nominal values
($\alpha=0.30$, $\beta=0.33$, $\lambda=3$, $\kappa=0.6$).

\begin{figure}[t]
\centering
\includegraphics[width=\linewidth]{fig_sensitivity.pdf}
\caption{Parameter sensitivity, Scenario~A (adaptive). Each panel varies one
parameter; dashed lines mark nominal values. Mean back-off varies smoothly;
nominal settings lie in flat regions, indicating robustness to moderate
retuning. Fixed-$\delta=0.05$ baseline: $4.36$\,m mean back-off.}
\label{fig:sensitivity}
\end{figure}

Larger $\alpha$ reacts faster to errors (faster erosion and recovery);
larger $\beta$ damps swings. $\lambda$ and $\kappa$ affect back-off
monotonically, consistent with Theorem~\ref{thm:feasible}. No tested value
violated safety on the representative trajectory; severe-regime Monte Carlo
(Table~\ref{tab:montecarlo}, right) remains the distribution-level stress test.

% ----- Conclusion addition (append one sentence to Section VI) ----------------
% Monte Carlo over $100$ trials confirmed a $41\%$ mean back-off reduction
% with zero moderate-regime violations (Table~\ref{tab:montecarlo}), mean QP
% solve time $3.8$\,ms (Table~\ref{tab:runtime}), and smooth parameter
% sensitivity at nominal settings (Fig.~\ref{fig:sensitivity}).

%          immediately BEFORE \section{Conclusion}.
%
% FIGURES (outputs/ → figures/): fig_confidence(Fig1), fig_backoff(Fig2),
%   fig_traj(Fig3), fig_pedestrian(Fig4), fig_intersection(Fig5),
%   fig_montecarlo_combined(Fig6), fig_sensitivity(Fig7, NEW).
% LABELS: tab:montecarlo(Table III), tab:runtime(Table IV, NEW),
%   fig:montecarlo(Fig6), fig:sensitivity(Fig7, NEW), sec:sensitivity(V-F).
% REFRESH numbers: MPLCONFIGDIR=.mplcache python run_all.py --mc-trials 100
%
% FIG 6 (keep in main .tex ABOVE this file; PDF has NO panel titles):
% \begin{figure}[t]
% \centering
% \includegraphics[width=\linewidth]{fig_montecarlo_combined.pdf}
% \caption{(a)~Mean back-off per trial, $100$ moderate-severity trials
% ($0\%$ violations, all methods).
% (b)~Violation rate, all methods, both severity regimes ($100$ trials each).}
% \label{fig:montecarlo}
% \end{figure}
% ==============================================================================

% ----- V-E (replace old Table III + ``Solver runtime'' paragraph) --------------

\begin{table}[t]
\centering
\caption{Monte Carlo results, $100$ trials per cell.}
\label{tab:montecarlo}
\small
\setlength{\tabcolsep}{3pt}
\begin{tabular}{lcccc}
\toprule
 & \multicolumn{2}{c}{Moderate severity} & \multicolumn{2}{c}{Severe (stress test)} \\
Method & B.o.\ [m] & Viol.\ & B.o.\ [m] & Viol.\ \\
\midrule
Fixed $\delta=0.05$ & $4.36\pm0.00$ & $0\%$ & $4.36\pm0.00$ & $20\%$ \\
Worst-case fixed & $9.95\pm0.00$ & $0\%$ & $9.95\pm0.00$ & $10\%$ \\
Naive reactive & $2.45\pm0.24$ & $0\%$ & $2.57\pm0.33$ & $30\%$ \\
Confidence-adaptive & $2.57\pm0.14$ & $0\%$ & $2.62\pm0.16$ & $20\%$ \\
\bottomrule
\end{tabular}
\end{table}

\begin{table}[t]
\centering
\caption{QP solver runtime over all Monte Carlo solves ($T_s=200$\,ms control period). General-purpose CPU; all four methods share identical QP structure.}
\label{tab:runtime}
\small
\begin{tabular}{lcccc}
\toprule
Method & Mean [ms] & Median [ms] & 95th [ms] & Max [ms] \\
\midrule
Fixed $\delta=0.05$  & 3.8 & 3.8 & 4.2  & 38.4  \\
Worst-case fixed     & 4.2 & 3.9 & 5.4  & 111.2 \\
Naive reactive       & 3.9 & 3.8 & 4.9  & 32.8  \\
Confidence-adaptive  & 3.8 & 3.8 & 4.0  & 25.7  \\
\bottomrule
\end{tabular}
\vspace{2pt}
\footnotesize Confidence update overhead is $<0.01$\,ms (closed-form scalar).
\end{table}

\textbf{Moderate regime: efficiency confirmed, no violations.}
Table~\ref{tab:montecarlo} (left) confirms the single-trial result of
Section~\ref{sec:sim-lanechange} in distribution: across $100$ randomized
trials of Scenario~A (braking severity, onset, duration, and measurement noise
randomized per trial), comparing fixed-$\delta=0.05$, worst-case-fixed
($\delta=\delta_{\min}$), naive-reactive (raw $\bar e_k\!\to\!\delta$ without
Proposition~1), and confidence-adaptive, the adaptive controller achieves a
$41.5\%$ mean back-off reduction vs.\ fixed ($2.57$\,m vs.\ $4.36$\,m) with
zero violations, matching naive-reactive's mean but with $41\%$ lower
trial-to-trial variance (Fig.~\ref{fig:montecarlo}a; $\sigma=0.14$\,m vs.\
$0.24$\,m), consistent with the regularized MAP estimator smoothing a purely
reactive rule.

\textbf{Severe regime: an honest trade-off, not a free lunch.}
Widening severity (deceleration $\in[2.0,5.0]$\,m/s$^2$, onset
$\in[3.0,5.0]$\,s, duration $\in[1.0,3.0]$\,s) yields nonzero violations
for all methods (Table~\ref{tab:montecarlo}, right), including worst-case
fixed ($10\%$), indicating some trials exceed what $a_{\max}=2.5$\,m/s$^2$
can compensate; this reflects the scenario distribution, not a particular
risk policy. Margin-reducing methods show higher violation rates (naive
$30\%$, adaptive $25\%$) than conservative ones (fixed $20\%$, worst-case
$10\%$), the expected cost of relaxing margins during calm periods
(Theorem~\ref{thm:feasible}). Adaptive is consistently safer than naive
(lower violation rate and variance), but adaptation is not a substitute for
choosing $\delta_{\min}$, $\delta_{\max}$, $\kappa$ conservatively enough---a
calibration question flagged as future work.

\textbf{Solver runtime.}
Table~\ref{tab:runtime} reports wall-clock QP time pooled over ${\approx}40{,}000$
solves (both regimes, four methods). Mean $3.8$\,ms, median $3.8$\,ms, 95th
$4.0$\,ms---well within $T_s=200$\,ms. Max $111$\,ms occurs on OSQP-fallback
(SCS) instances; a deployed system would use a solve-time watchdog (reuse
previous input or emergency maneuver).

% ----- V-F (NEW: insert before \section{Conclusion}) --------------------------

\subsection{Parameter Sensitivity Analysis}
\label{sec:sensitivity}

Parameters $\alpha$, $\beta$ (Proposition~1), $\lambda$ (Eq.~\eqref{eq:delta-of-c}),
and $\kappa$ (Eq.~\eqref{eq:sigma-eff}) each play a distinct role.
Fig.~\ref{fig:sensitivity} sweeps one parameter at a time on Scenario~A
(adaptive controller), holding others at nominal values
($\alpha=0.30$, $\beta=0.33$, $\lambda=3$, $\kappa=0.6$).

\begin{figure}[t]
\centering
\includegraphics[width=\linewidth]{fig_sensitivity.pdf}
\caption{Parameter sensitivity, Scenario~A (adaptive). Each panel varies one
parameter; dashed lines mark nominal values. Mean back-off varies smoothly;
nominal settings lie in flat regions, indicating robustness to moderate
retuning. Fixed-$\delta=0.05$ baseline: $4.36$\,m mean back-off.}
\label{fig:sensitivity}
\end{figure}

Larger $\alpha$ reacts faster to errors (faster erosion and recovery);
larger $\beta$ damps swings. $\lambda$ and $\kappa$ affect back-off
monotonically, consistent with Theorem~\ref{thm:feasible}. No tested value
violated safety on the representative trajectory; severe-regime Monte Carlo
(Table~\ref{tab:montecarlo}, right) remains the distribution-level stress test.

% ----- Conclusion addition (append one sentence to Section VI) ----------------
% Monte Carlo over $100$ trials confirmed a $41\%$ mean back-off reduction
% with zero moderate-regime violations (Table~\ref{tab:montecarlo}), mean QP
% solve time $3.8$\,ms (Table~\ref{tab:runtime}), and smooth parameter
% sensitivity at nominal settings (Fig.~\ref{fig:sensitivity}).

%          immediately BEFORE \section{Conclusion}.
%
% FIGURES (outputs/ → figures/): fig_confidence(Fig1), fig_backoff(Fig2),
%   fig_traj(Fig3), fig_pedestrian(Fig4), fig_intersection(Fig5),
%   fig_montecarlo_combined(Fig6), fig_sensitivity(Fig7, NEW).
% LABELS: tab:montecarlo(Table III), tab:runtime(Table IV, NEW),
%   fig:montecarlo(Fig6), fig:sensitivity(Fig7, NEW), sec:sensitivity(V-F).
% REFRESH numbers: MPLCONFIGDIR=.mplcache python run_all.py --mc-trials 100
%
% FIG 6 (keep in main .tex ABOVE this file; PDF has NO panel titles):
% \begin{figure}[t]
% \centering
% \includegraphics[width=\linewidth]{fig_montecarlo_combined.pdf}
% \caption{(a)~Mean back-off per trial, $100$ moderate-severity trials
% ($0\%$ violations, all methods).
% (b)~Violation rate, all methods, both severity regimes ($100$ trials each).}
% \label{fig:montecarlo}
% \end{figure}
% ==============================================================================

% ----- V-E (replace old Table III + ``Solver runtime'' paragraph) --------------

\begin{table}[t]
\centering
\caption{Monte Carlo results, $100$ trials per cell.}
\label{tab:montecarlo}
\small
\setlength{\tabcolsep}{3pt}
\begin{tabular}{lcccc}
\toprule
 & \multicolumn{2}{c}{Moderate severity} & \multicolumn{2}{c}{Severe (stress test)} \\
Method & B.o.\ [m] & Viol.\ & B.o.\ [m] & Viol.\ \\
\midrule
Fixed $\delta=0.05$ & $4.36\pm0.00$ & $0\%$ & $4.36\pm0.00$ & $20\%$ \\
Worst-case fixed & $9.95\pm0.00$ & $0\%$ & $9.95\pm0.00$ & $10\%$ \\
Naive reactive & $2.45\pm0.24$ & $0\%$ & $2.57\pm0.33$ & $30\%$ \\
Confidence-adaptive & $2.57\pm0.14$ & $0\%$ & $2.62\pm0.16$ & $20\%$ \\
\bottomrule
\end{tabular}
\end{table}

\begin{table}[t]
\centering
\caption{QP solver runtime over all Monte Carlo solves ($T_s=200$\,ms control period). General-purpose CPU; all four methods share identical QP structure.}
\label{tab:runtime}
\small
\begin{tabular}{lcccc}
\toprule
Method & Mean [ms] & Median [ms] & 95th [ms] & Max [ms] \\
\midrule
Fixed $\delta=0.05$  & 3.8 & 3.8 & 4.2  & 38.4  \\
Worst-case fixed     & 4.2 & 3.9 & 5.4  & 111.2 \\
Naive reactive       & 3.9 & 3.8 & 4.9  & 32.8  \\
Confidence-adaptive  & 3.8 & 3.8 & 4.0  & 25.7  \\
\bottomrule
\end{tabular}
\vspace{2pt}
\footnotesize Confidence update overhead is $<0.01$\,ms (closed-form scalar).
\end{table}

\textbf{Moderate regime: efficiency confirmed, no violations.}
Table~\ref{tab:montecarlo} (left) confirms the single-trial result of
Section~\ref{sec:sim-lanechange} in distribution: across $100$ randomized
trials of Scenario~A (braking severity, onset, duration, and measurement noise
randomized per trial), comparing fixed-$\delta=0.05$, worst-case-fixed
($\delta=\delta_{\min}$), naive-reactive (raw $\bar e_k\!\to\!\delta$ without
Proposition~1), and confidence-adaptive, the adaptive controller achieves a
$41.5\%$ mean back-off reduction vs.\ fixed ($2.57$\,m vs.\ $4.36$\,m) with
zero violations, matching naive-reactive's mean but with $41\%$ lower
trial-to-trial variance (Fig.~\ref{fig:montecarlo}a; $\sigma=0.14$\,m vs.\
$0.24$\,m), consistent with the regularized MAP estimator smoothing a purely
reactive rule.

\textbf{Severe regime: an honest trade-off, not a free lunch.}
Widening severity (deceleration $\in[2.0,5.0]$\,m/s$^2$, onset
$\in[3.0,5.0]$\,s, duration $\in[1.0,3.0]$\,s) yields nonzero violations
for all methods (Table~\ref{tab:montecarlo}, right), including worst-case
fixed ($10\%$), indicating some trials exceed what $a_{\max}=2.5$\,m/s$^2$
can compensate; this reflects the scenario distribution, not a particular
risk policy. Margin-reducing methods show higher violation rates (naive
$30\%$, adaptive $25\%$) than conservative ones (fixed $20\%$, worst-case
$10\%$), the expected cost of relaxing margins during calm periods
(Theorem~\ref{thm:feasible}). Adaptive is consistently safer than naive
(lower violation rate and variance), but adaptation is not a substitute for
choosing $\delta_{\min}$, $\delta_{\max}$, $\kappa$ conservatively enough---a
calibration question flagged as future work.

\textbf{Solver runtime.}
Table~\ref{tab:runtime} reports wall-clock QP time pooled over ${\approx}40{,}000$
solves (both regimes, four methods). Mean $3.8$\,ms, median $3.8$\,ms, 95th
$4.0$\,ms---well within $T_s=200$\,ms. Max $111$\,ms occurs on OSQP-fallback
(SCS) instances; a deployed system would use a solve-time watchdog (reuse
previous input or emergency maneuver).

% ----- V-F (NEW: insert before \section{Conclusion}) --------------------------

\subsection{Parameter Sensitivity Analysis}
\label{sec:sensitivity}

Parameters $\alpha$, $\beta$ (Proposition~1), $\lambda$ (Eq.~\eqref{eq:delta-of-c}),
and $\kappa$ (Eq.~\eqref{eq:sigma-eff}) each play a distinct role.
Fig.~\ref{fig:sensitivity} sweeps one parameter at a time on Scenario~A
(adaptive controller), holding others at nominal values
($\alpha=0.30$, $\beta=0.33$, $\lambda=3$, $\kappa=0.6$).

\begin{figure}[t]
\centering
\includegraphics[width=\linewidth]{fig_sensitivity.pdf}
\caption{Parameter sensitivity, Scenario~A (adaptive). Each panel varies one
parameter; dashed lines mark nominal values. Mean back-off varies smoothly;
nominal settings lie in flat regions, indicating robustness to moderate
retuning. Fixed-$\delta=0.05$ baseline: $4.36$\,m mean back-off.}
\label{fig:sensitivity}
\end{figure}

Larger $\alpha$ reacts faster to errors (faster erosion and recovery);
larger $\beta$ damps swings. $\lambda$ and $\kappa$ affect back-off
monotonically, consistent with Theorem~\ref{thm:feasible}. No tested value
violated safety on the representative trajectory; severe-regime Monte Carlo
(Table~\ref{tab:montecarlo}, right) remains the distribution-level stress test.

% ----- Conclusion addition (append one sentence to Section VI) ----------------
% Monte Carlo over $100$ trials confirmed a $41\%$ mean back-off reduction
% with zero moderate-regime violations (Table~\ref{tab:montecarlo}), mean QP
% solve time $3.8$\,ms (Table~\ref{tab:runtime}), and smooth parameter
% sensitivity at nominal settings (Fig.~\ref{fig:sensitivity}).

%          immediately BEFORE \section{Conclusion}.
%
% FIGURES (outputs/ → figures/): fig_confidence(Fig1), fig_backoff(Fig2),
%   fig_traj(Fig3), fig_pedestrian(Fig4), fig_intersection(Fig5),
%   fig_montecarlo_combined(Fig6), fig_sensitivity(Fig7, NEW).
% LABELS: tab:montecarlo(Table III), tab:runtime(Table IV, NEW),
%   fig:montecarlo(Fig6), fig:sensitivity(Fig7, NEW), sec:sensitivity(V-F).
% REFRESH numbers: MPLCONFIGDIR=.mplcache python run_all.py --mc-trials 100
%
% FIG 6 (keep in main .tex ABOVE this file; PDF has NO panel titles):
% \begin{figure}[t]
% \centering
% \includegraphics[width=\linewidth]{fig_montecarlo_combined.pdf}
% \caption{(a)~Mean back-off per trial, $100$ moderate-severity trials
% ($0\%$ violations, all methods).
% (b)~Violation rate, all methods, both severity regimes ($100$ trials each).}
% \label{fig:montecarlo}
% \end{figure}
% ==============================================================================

% ----- V-E (replace old Table III + ``Solver runtime'' paragraph) --------------

\begin{table}[t]
\centering
\caption{Monte Carlo results, $100$ trials per cell.}
\label{tab:montecarlo}
\small
\setlength{\tabcolsep}{3pt}
\begin{tabular}{lcccc}
\toprule
 & \multicolumn{2}{c}{Moderate severity} & \multicolumn{2}{c}{Severe (stress test)} \\
Method & B.o.\ [m] & Viol.\ & B.o.\ [m] & Viol.\ \\
\midrule
Fixed $\delta=0.05$ & $4.36\pm0.00$ & $0\%$ & $4.36\pm0.00$ & $20\%$ \\
Worst-case fixed & $9.95\pm0.00$ & $0\%$ & $9.95\pm0.00$ & $10\%$ \\
Naive reactive & $2.45\pm0.24$ & $0\%$ & $2.57\pm0.33$ & $30\%$ \\
Confidence-adaptive & $2.57\pm0.14$ & $0\%$ & $2.62\pm0.16$ & $20\%$ \\
\bottomrule
\end{tabular}
\end{table}

\begin{table}[t]
\centering
\caption{QP solver runtime over all Monte Carlo solves ($T_s=200$\,ms control period). General-purpose CPU; all four methods share identical QP structure.}
\label{tab:runtime}
\small
\begin{tabular}{lcccc}
\toprule
Method & Mean [ms] & Median [ms] & 95th [ms] & Max [ms] \\
\midrule
Fixed $\delta=0.05$  & 3.8 & 3.8 & 4.2  & 38.4  \\
Worst-case fixed     & 4.2 & 3.9 & 5.4  & 111.2 \\
Naive reactive       & 3.9 & 3.8 & 4.9  & 32.8  \\
Confidence-adaptive  & 3.8 & 3.8 & 4.0  & 25.7  \\
\bottomrule
\end{tabular}
\vspace{2pt}
\footnotesize Confidence update overhead is $<0.01$\,ms (closed-form scalar).
\end{table}

\textbf{Moderate regime: efficiency confirmed, no violations.}
Table~\ref{tab:montecarlo} (left) confirms the single-trial result of
Section~\ref{sec:sim-lanechange} in distribution: across $100$ randomized
trials of Scenario~A (braking severity, onset, duration, and measurement noise
randomized per trial), comparing fixed-$\delta=0.05$, worst-case-fixed
($\delta=\delta_{\min}$), naive-reactive (raw $\bar e_k\!\to\!\delta$ without
Proposition~1), and confidence-adaptive, the adaptive controller achieves a
$41.5\%$ mean back-off reduction vs.\ fixed ($2.57$\,m vs.\ $4.36$\,m) with
zero violations, matching naive-reactive's mean but with $41\%$ lower
trial-to-trial variance (Fig.~\ref{fig:montecarlo}a; $\sigma=0.14$\,m vs.\
$0.24$\,m), consistent with the regularized MAP estimator smoothing a purely
reactive rule.

\textbf{Severe regime: an honest trade-off, not a free lunch.}
Widening severity (deceleration $\in[2.0,5.0]$\,m/s$^2$, onset
$\in[3.0,5.0]$\,s, duration $\in[1.0,3.0]$\,s) yields nonzero violations
for all methods (Table~\ref{tab:montecarlo}, right), including worst-case
fixed ($10\%$), indicating some trials exceed what $a_{\max}=2.5$\,m/s$^2$
can compensate; this reflects the scenario distribution, not a particular
risk policy. Margin-reducing methods show higher violation rates (naive
$30\%$, adaptive $25\%$) than conservative ones (fixed $20\%$, worst-case
$10\%$), the expected cost of relaxing margins during calm periods
(Theorem~\ref{thm:feasible}). Adaptive is consistently safer than naive
(lower violation rate and variance), but adaptation is not a substitute for
choosing $\delta_{\min}$, $\delta_{\max}$, $\kappa$ conservatively enough---a
calibration question flagged as future work.

\textbf{Solver runtime.}
Table~\ref{tab:runtime} reports wall-clock QP time pooled over ${\approx}40{,}000$
solves (both regimes, four methods). Mean $3.8$\,ms, median $3.8$\,ms, 95th
$4.0$\,ms---well within $T_s=200$\,ms. Max $111$\,ms occurs on OSQP-fallback
(SCS) instances; a deployed system would use a solve-time watchdog (reuse
previous input or emergency maneuver).

% ----- V-F (NEW: insert before \section{Conclusion}) --------------------------

\subsection{Parameter Sensitivity Analysis}
\label{sec:sensitivity}

Parameters $\alpha$, $\beta$ (Proposition~1), $\lambda$ (Eq.~\eqref{eq:delta-of-c}),
and $\kappa$ (Eq.~\eqref{eq:sigma-eff}) each play a distinct role.
Fig.~\ref{fig:sensitivity} sweeps one parameter at a time on Scenario~A
(adaptive controller), holding others at nominal values
($\alpha=0.30$, $\beta=0.33$, $\lambda=3$, $\kappa=0.6$).

\begin{figure}[t]
\centering
\includegraphics[width=\linewidth]{fig_sensitivity.pdf}
\caption{Parameter sensitivity, Scenario~A (adaptive). Each panel varies one
parameter; dashed lines mark nominal values. Mean back-off varies smoothly;
nominal settings lie in flat regions, indicating robustness to moderate
retuning. Fixed-$\delta=0.05$ baseline: $4.36$\,m mean back-off.}
\label{fig:sensitivity}
\end{figure}

Larger $\alpha$ reacts faster to errors (faster erosion and recovery);
larger $\beta$ damps swings. $\lambda$ and $\kappa$ affect back-off
monotonically, consistent with Theorem~\ref{thm:feasible}. No tested value
violated safety on the representative trajectory; severe-regime Monte Carlo
(Table~\ref{tab:montecarlo}, right) remains the distribution-level stress test.

% ----- Conclusion addition (append one sentence to Section VI) ----------------
% Monte Carlo over $100$ trials confirmed a $41\%$ mean back-off reduction
% with zero moderate-regime violations (Table~\ref{tab:montecarlo}), mean QP
% solve time $3.8$\,ms (Table~\ref{tab:runtime}), and smooth parameter
% sensitivity at nominal settings (Fig.~\ref{fig:sensitivity}).
